\section{\texorpdfstring{$N=1$}{N=1} Boundary-Compatibility Conditions}
\label{sec:n1-boundary-compatibility-conditions}

Let a chambered Hall-descent datum supply the compact comparison data of
Definition~\ref{def:compact-igusa-realization-datum}, so that the
criterion of Theorem~\ref{thm:primitive-recognition} applies.
This appendix uses the same smooth compact-realization standing datum:
\(S\) is a smooth complex K3 surface and \(E\) is a smooth elliptic curve.
Singular K3 limits, nodal or cuspidal elliptic limits, imprimitive
multiple-cover contributions, the \(s=0\) compact Hall cusp, and
multi-component charge closures require separate boundary-degeneration
certificates.
Its \(N=1\) decategorified boundary would then be constrained by the
Igusa datum
\[
\begin{aligned}
\Gamma_{\mathrm{eff}}
= {}&
\{(n,l,m)\in\Z^3\mid m>0,\ n\ge0,\ l\in\Z\}\\
&{}\cup\{(n,l,0)\in\Z^3\mid n>0,\ l\in\Z\}\\
&{}\cup\{(0,l,0)\in\Z^3\mid l<0\}
\subset \Z^3 .
\end{aligned}
\]

\[
\begin{array}{c|c}
\mathrm{comparison\ Hall\ descent} & \mathrm{Igusa\ boundary}\\ \hline
\mathrm{Gram\ index\ lattice} & \Gamma_{\mathrm{ind}}=\Z^3\\
\mathrm{sector\ completion} & \Gamma_{\mathrm{eff}}\\
\mathrm{active\ denominator\ support} &
\Gamma_{\mathrm{act}}=\{\gamma\in\Gamma_{\mathrm{eff}}\mid f(nm,l)\ne0\}\\
\mathrm{Lorentzian\ degree} &
\alpha(n,l,m)=2nf_2-lf_3+2mf_{-2}\\
\mathrm{orientation\ character} &
\hbox{required geometric monodromy comparison }
\epsilon_o|_{W^{(2)}}=\nu_{\Delta_5}|_{W^{(2)}}\\
\mathrm{automorphic\ boundary} &
\operatorname{Borch}_{\theta}(\phi_{0,1})=\Delta_5\\
\mathrm{BKM\ denominator} &
\operatorname{den}(\autcor)=64^{-1}\Delta_5(2Z)\\
\mathrm{denominator\ algebra} &
\mathcal A_{\mathrm{den}}=\mathfrak g_{\Delta_5}\\
\mathrm{formal\ current\ envelope} &
\mathsf{Den}_{\Delta_5,E}
=U_E^{\mathrm{ch}}(\mathfrak g_{\Delta_5}\otimes\mathcal O_E)\\
\mathrm{scalar\ square} &
Z^X_{\mathrm{OP}}=Z^X_{\mathrm{DT}}=-4096\Delta_5^{-2}
\quad\text{(proved for the smooth primitive nonzero \(K3\)-class branch).}
\end{array}
\]

The degree map satisfies
\[
(\alpha(\gamma),\alpha(\eta))
=-2\langle\gamma,\eta\rangle_{\mathrm{ind}},
\qquad
(\alpha(n,l,m),\alpha(n,l,m))=-2(4nm-l^2).
\]
The real walls are
\[
\delta_1=2f_2-f_3,\qquad
\delta_2=2f_{-2}-f_3,\qquad
\delta_3=f_3.
\]
The Weyl vector is
\[
\rho=f_2-\frac12 f_3+f_{-2},\qquad
(\rho,\delta_i)=-1.
\]
These equations give the denominator polarisation of the same
Gram-index lattice whose cusp polarisation gives the Fock determinant.
The cone is generated by the active support.  Triples in
\(\Gamma_{\mathrm{eff}}\) with zero Jacobi coefficient fix ordering and
contribute no visible product factors.  The product alone does not
exclude invisible zero-superdimension root spaces.

The comparison is the \(N=1\) normalization
\[
\operatorname{Borch}_{\theta}(\phi_{0,1})=\Delta_5,\qquad
\operatorname{den}(\autcor)=64^{-1}\Delta_5(2Z),\qquad
\kappa_{\mathrm{BKM}}(\autcor)=\mathrm{wt}(\Delta_5)=5.
\]
It contains no denominator theorem beyond the \(N=1\) theorem and no row
map
\[
\Phi^{\mathrm{cmp}}_j\longleftrightarrow F_i.
\]
The protected scalar has the form
\[
Z^X_{\square}=C_{\square}\Delta_5^{-2}.
\]
The Oberdieck--Pixton and reduced Hilbert-scheme DT branches have
constant \(-4096\) in the quotient-by-\(E\) normalization of
Theorem~\ref{thm:op-square}.  The chiral/denominator half is
\(\Delta_5\).

A compact Hall--Drinfeld double would realize a protected bracket from
the compactification only when the comparison package of
Definition~\ref{def:compact-igusa-realization-datum} and
Theorem~\ref{thm:primitive-recognition} is supplied:
finite-first oriented Hall descent, protected integration, charge
additive Hall primitives, Borcherds relations inside the Hall
super-commutator, and the completed Hopf pairing with orientation
character required to restrict to \(\nu_{\Delta_5}\).  Homotopy descent
must produce the positive half, its completion data, and the Hall-side
domain on which the comparison character is defined.  This orientation
character is not produced by rescaling \(\Delta_5\); it must come from
monodromy of the chosen square root of the Hall determinant line.  The
Igusa boundary identifies the denominator algebra
\(\mathfrak g_{\Delta_5}\) and its protected decategorified
multiplicities before those comparison data are proved.

For the Cl\'ery--Gritsenko rows \(F_5,F_6,F_7,F_8\), this table contains
no nonabelian chamber, no simple real roots, no reflection character,
and no Weyl-sum side.  Rows \(F_5,F_6,F_7,F_8\) use the separate
Govindarajan CHL/WKB denominator theorem when a nonabelian bracket is
needed.  For \(F_7\) this is a theorem on the \(\mu_4\)-multiplier cover
determined by the Cl\'ery--Gritsenko order-four multiplier.  None of these
denominator rows supplies a comparison row map or a CY/DT/chiral
boundary host.
